\chapter{Instruments delta-one}\label{ch:m1:delta-one-instruments}

Cinq investisseurs veulent la même chose~: le rendement d'un indice actions
sur cent millions d'euros pendant un an. Le premier achète les actions, le deuxième un fonds, le
troisième des contrats à terme, le quatrième signe un swap avec une banque, le cinquième ouvre un
\omterm{def:m1:delta-one-instruments:cfd}{contrat sur différence} auprès d'un courtier. À la fin de l'année, l'indice a
rapporté ce qu'il a rapporté, et les cinq ont gagné des montants différents.
La voie la plus coûteuse a coûté plus de cinq fois la moins chère, et la voie la moins chère
dépend de qui est l'investisseur~: s'il dispose du cash, de l'impôt qu'il paie
sur les dividendes, de la durée de sa détention. Aucun des cinq instruments
ne contient une option, une opinion ou un modèle. Ils diffèrent par le financement, par la fiscalité et par
la plomberie, c'est-à-dire par tout ce dont les seize premiers chapitres
traitaient.

\section{Ce que signifie delta-one}

\begin{definition}[Delta-one]\label{def:m1:delta-one-instruments:deltaone}
Un instrument est \emph{delta-one}\index{delta-one} sur un sous-jacent lorsque
sa valeur varie d'une unité pour une variation d'une unité de la valeur du
sous-jacent~: il porte l'exposition sans optionnalité. Les actions, les fonds,
les contrats à terme et les contrats à terme de gré à gré, les \omterm{def:m1:financing:trs}{total return swaps}, les \omterm{def:m1:delta-one-instruments:cfd}{contrats sur différence},
les \omterm{def:m1:delta-one-instruments:dr}{certificats de dépôt} et les certificats qui répliquent un actif un pour un sont
delta-one. Un \emph{desk delta-one} fixe les prix, couvre et arbitre les
écarts qui les séparent.
\end{definition}

Puisque tous les véhicules offrent la même exposition, seuls leurs coûts peuvent différer.
Il existe cinq types de coûts, et chaque véhicule est un vecteur de cinq nombres.

\begin{method}[Les cinq coûts d'un véhicule]\label{met:m1:delta-one-instruments:five}
\begin{enumerate}
\item \textbf{Négociation}~: commissions, fourchette et impact pour entrer et sortir,
et pour faire le roll des instruments qui arrivent à échéance.
\item \textbf{Taxe sur les transactions} à l'achat, lorsqu'elle existe.
\item \textbf{Frais courants}~: les \omterm{def:m1:the-buy-side:fees}{frais de gestion} d'un fonds.
\item \textbf{Financement}~: le \omterm{def:m1:delta-one-instruments:swap}{spread de financement}, soit la marge au-dessus du taux de référence payé sur l'argent
emprunté, explicitement (un prêt sur marge) ou dans le prix (un contrat à terme,
un swap).
\item \textbf{Dividendes}~: la fraction perdue à la \omterm{def:m1:delta-one-instruments:wht}{retenue à la source}, qui dépend
du détenteur \emph{et} du véhicule.
\end{enumerate}
Exprimer chacun en points de base de notionnel sur la durée de détention visée,
puis additionner.
\end{method}

\section{Swaps, contrats sur différence et spread de financement}

\begin{definition}[Equity swap et spread de financement]\label{def:m1:delta-one-instruments:swap}
Un \emph{equity swap}\index{equity swap} est un \omterm{def:m1:financing:trs}{total return swap}
(\cref{def:m1:financing:trs}) sur une action, un panier ou un indice, avec des
recalages périodiques auxquels la performance courue est versée et le notionnel est
recalé sur le cours courant. Son \emph{spread de financement}\index{spread de financement} est la
marge au-dessus du taux de référence dans la jambe de financement~: le prix du
bilan du dealer, et le chiffre sur lequel les swaps se font concurrence.
\end{definition}

\begin{omfigure}
\begin{tikzpicture}[font=\small,
  box/.style={draw=omDef,rounded corners=2pt,minimum width=2.2cm,minimum height=1.1cm,align=center,fill=omDef!4},
  flow/.style={-{Stealth[length=2mm]},thick}]
\node[box] (c) at (0,0) {Client};
\node[box] (d) at (7.0,0) {Dealer};
\node[box,draw=omExo,fill=omExo!6] (m) at (12.2,0) {Marché};
\draw[flow,omDef] ([yshift=9pt]d.west) -- node[above,font=\footnotesize,align=center]{rendement du cours\\$+$ une part $\eta$ des dividendes} ([yshift=9pt]c.east);
\draw[flow,omThm] ([yshift=-9pt]c.east) -- node[below,font=\footnotesize,align=center]{taux de référence\\$+$ spread de financement} ([yshift=-9pt]d.west);
\draw[flow,omExo] (d.east) -- node[above,font=\footnotesize]{achète les actions} node[below,font=\footnotesize]{en couverture} (m.west);
\end{tikzpicture}

\omcaption{Un \omterm{def:m1:delta-one-instruments:swap}{equity swap}. Le dealer est à plat~: il possède ce qu'il doit. Son
profit est le \omterm{def:m1:delta-one-instruments:swap}{spread de financement}, moins son propre coût de financement des actions, plus
tout ce que les actions rapportent entre ses mains (frais de \omterm{def:m1:financing:seclending}{prêt de titres}, meilleure
position fiscale sur les dividendes) qu'il ne reverse pas.}\label{fig:m1:delta-one-instruments:swap}
\end{omfigure}

\begin{definition}[Contrat sur différence]\label{def:m1:delta-one-instruments:cfd}
Un \emph{contrat sur différence}\index{contrat sur différence} (CFD) est un
\omterm{def:m1:delta-one-instruments:swap}{equity swap} à durée indéterminée proposé par un courtier à ses clients, en général des particuliers,
sur marge~: le client reçoit la variation du cours, se voit créditer ou débiter des
dividendes, et paie un financement quotidien sur le notionnel complet tant que
la position reste ouverte.
\end{definition}

Un swap et un CFD sont le même contrat à deux niveaux de prix. L'institution
négocie un \omterm{def:m1:delta-one-instruments:swap}{spread de financement} de quelques dizaines de points de base auprès de plusieurs dealers~; le
client particulier prend le taux du courtier, souvent de quelques \emph{pour cent} au-dessus
du taux de référence. Les deux existent pour les mêmes deux raisons~: l'effet de levier sans
prêt sur marge, et le fait que le client ne détient jamais l'action.

\begin{definition}[Retenue à la source et optimisation des dividendes]\label{def:m1:delta-one-instruments:wht}
Une \emph{retenue à la source}\index{retenue à la source} est un impôt prélevé à la source
sur les dividendes versés à un détenteur non résident, à un taux légal que peut réduire une
convention fiscale. L'\emph{optimisation des dividendes}\index{optimisation des dividendes} désigne tout
arrangement dans lequel l'exposition est détenue, à la date du dividende, par la
partie soumise au coût fiscal le plus faible, qui reverse une partie de l'économie au
détenteur économique, typiquement par le biais de la part $\eta$ du dividende versée dans
un swap.
\end{definition}

\begin{dated}{2026-09}{Deux taxes qui façonnent les véhicules}\label{dat:m1:delta-one-instruments:tax}
\textbf{Royaume-Uni.} L'achat électronique d'actions britanniques coûte 0.5~\% de Stamp
Duty Reserve Tax~; le transfert d'actions vers certains programmes de \omterm{def:m1:delta-one-instruments:dr}{certificats de dépôt} ou
services de compensation coûte 1.5~\%~; les actions étrangères achetées hors du Royaume-Uni ne sont
normalement pas taxées. Les intermédiaires agréés, comme les \omterm{def:m1:what-a-trading-firm-does:market-maker}{teneurs de marché}, sont
exonérés de la taxe sur leurs achats, ce qui permet à un dealer de couvrir un
swap ou un CFD, sur lesquels aucune action ne change de mains, sans la payer.

\textbf{États-Unis.} La section 871(m) du code fiscal traite les paiements
\emph{dividend-equivalent} sur des dérivés portant sur des actions américaines comme des dividendes de source américaine,
soumis à la même retenue que les dividendes eux-mêmes. Elle
s'applique aux contrats \omterm{def:m1:delta-one-instruments:deltaone}{delta-one}~; selon la Notice 2024-44, les contrats qui ne sont pas
\omterm{def:m1:delta-one-instruments:deltaone}{delta-one} et qui sont émis avant le 1\up{er}~janvier 2027 restent hors de son champ. Pour un
détenteur non américain, un swap sur une action américaine n'est plus avantageux que l'action elle-même.
\end{dated}

\section{Contrats à terme comme produits delta-one}

Un contrat à terme (partie III) est l'enveloppe la plus standardisée~: pas de contrepartie avec laquelle négocier,
une marge de quelques pour cent et le financement inclus dans le prix.
Si l'indice vaut $S$, que le taux de référence est $r$ et que le rendement des dividendes est $q$,
le contrat à terme d'échéance $T$ est au juste prix lorsqu'il vaut $S(1+(r-q)T)$~; le taux qui rend le
prix de \emph{marché} juste est le \emph{taux de financement implicite} du contrat à terme, et
son excédent sur le taux de référence est le \omterm{def:m1:delta-one-instruments:swap}{spread de financement} du marché des contrats à terme.
Il est fixé par l'offre et la demande de levier~: lorsque beaucoup d'investisseurs veulent être
acheteurs via des contrats à terme et que les bilans des dealers sont rares, comme en
fin d'année, le spread monte. Le détenteur paie aussi un roll, quatre fois par an,
dont le coût est ce spread pour le trimestre suivant, auquel s'ajoute le coût de transaction
du calendar spread (\cref{ch:m1:basis-roll-and-delivery}).

\begin{example}[Lire le spread à partir d'un prix]\label{ex:m1:delta-one-instruments:implied}
Indice à 6\,000, $r = 4.2\%$, $q = 1.3\%$, trois mois~: fair value $6\,000
\times (1 + 0.029 \times 0.25) = 6\,043.5$. Le contrat à terme s'échange à 6\,048. Son
taux de financement implicite est $(6\,048/6\,000 - 1)/0.25 + 1.3\% = 4.5\%$~: un
spread de 30 points de base. Un investisseur disposant de liquidités, qui achète le contrat à terme
et place ses liquidités au taux de référence, paie 30 points de base par an pour
l'exposition~; un investisseur qui emprunterait autrement auprès d'un \omterm{def:m1:financing:pb}{prime broker} au taux
de référence majoré de 50 points de base \emph{économise} 20 points de base.
\end{example}

\section{Certificats de dépôt et doubles cotations}

\begin{definition}[Certificat de dépôt]\label{def:m1:delta-one-instruments:dr}
Un \emph{certificat de dépôt}\index{certificat de dépôt} est un titre émis par une banque
dépositaire qui représente un nombre fixe (le \emph{ratio}) d'actions
d'une société étrangère conservées en dépôt sur le marché d'origine. Il se négocie et se règle
dans la devise et le système de la place d'accueil~; les American depositary
receipts (ADRs) sont la forme américaine. Les certificats peuvent être \emph{émis} contre un
dépôt d'actions ordinaires et \emph{annulés} pour la livraison des actions,
moyennant des frais par certificat.
\end{definition}

\begin{definition}[Double cotation]\label{def:m1:delta-one-instruments:dual}
Une société dispose d'une \emph{double cotation}\index{double cotation} lorsque la même classe
d'actions est cotée sur des bourses de deux pays et peut être transférée entre
les deux registres. Contrairement à un certificat, il n'existe pas de titre intermédiaire,
seulement deux lieux où la même action se négocie.
\end{definition}

\begin{dated}{2012-08}{Ce que dit le bulletin du régulateur sur les coûts des ADR}\label{dat:m1:delta-one-instruments:adr}
Les banques dépositaires facturent aux détenteurs des frais de garde (depositary services fee),
généralement retenus sur le dividende brut~; ces frais sont en général calculés par
ADR, l'exemple du bulletin étant de 20 à 50~\$ pour 1\,000 ADRs. Le barème des frais
de chaque programme figure dans son formulaire F-6. Les programmes parrainés ont
trois niveaux~: le niveau 1 se négocie de gré à gré, le niveau 2 est coté en bourse,
le niveau 3 peut lever des capitaux. Les ADR non parrainés sont mis en place sans
la coopération de la société.
\end{dated}

\begin{proposition}[La bande de conversion d'un certificat]\label{prop:m1:delta-one-instruments:band}
Soit $P^* = \rho\,S\,X$ le prix de parité d'un certificat de ratio $\rho$,
de prix ordinaire $S$ et de taux de change $X$. Avec des coûts de transaction $c$ (les deux
jambes et l'opération de change, exprimée en fraction), des frais d'émission et d'annulation $f_i$ et
$f_c$ par certificat, et une taxe $\tau$ sur le dépôt d'actions, la conversion n'est
rentable que lorsque la prime du certificat par rapport à la parité se situe hors de
\[
\Bigl[-\bigl(c + f_c/P^*\bigr),\; c + \tau + f_i/P^*\Bigr].
\]
Des frais fixes en centimes élargissent la bande des certificats à bas prix, et une
taxe sur le dépôt la rend asymétrique.
\end{proposition}
\begin{proof}
Au-dessus de la parité~: acheter $\rho$ actions ordinaires, payer $\tau$ et $f_i$, recevoir un
certificat, le vendre~; le profit par certificat est $P - P^*(1+c+\tau) - f_i$. En dessous~: acheter le
certificat, payer $f_c$, recevoir et vendre les actions ordinaires.
\end{proof}

\begin{omfigure}
\begin{tikzpicture}
\begin{axis}[width=10cm,height=4.8cm,
  xlabel={prix du certificat (\$)},ylabel={prime par rapport à la parité (pb)},
  tick label style={font=\small},label style={font=\small},
  xmin=0,xmax=105,ymin=-120,ymax=120,grid=major,grid style={gray!25},table/col sep=comma]
\addplot[omDef,very thick,mark=*,mark size=1.5pt] table[x=receipt_price,y=upper_bp]{figdata/markets-1/17-delta-one-instruments/band.csv};
\addplot[omDef,very thick,mark=*,mark size=1.5pt] table[x=receipt_price,y=lower_bp]{figdata/markets-1/17-delta-one-instruments/band.csv};
\node[font=\footnotesize,anchor=west,align=left] (in) at (axis cs:64,-72) {à l'intérieur~:\\rien n'est rentable};
\draw[-{Stealth[length=1.6mm]}] (in.north) -- (axis cs:80,-2);
\node[font=\footnotesize,anchor=west] at (axis cs:32,72) {émettre des certificats};
\node[font=\footnotesize,anchor=west] at (axis cs:22,-72) {annuler des certificats};
\end{axis}
\end{tikzpicture}

\omcaption{La bande de conversion pour des frais de 5 cents par certificat dans chaque sens et
des coûts de transaction de 6 points de base, sans taxe sur le dépôt. À 5~\$, les frais seuls
représentent 100 points de base. Données~:
\cref{prop:m1:delta-one-instruments:band}.}\label{fig:m1:delta-one-instruments:band}
\end{omfigure}

La parité n'est observable que tant que les deux marchés sont ouverts. Lorsque le marché
d'origine est fermé, le certificat est le seul cours disponible, et sa «~prime~» par rapport
à la clôture périmée du marché d'origine constitue une information, comme pour les ETF présentés dans
la \cref{prop:m1:exchange-traded-funds:stale}. Certaines paires ne convergent jamais du tout~:
lorsque la conversion est restreinte par des contrôles de capitaux ou des limites de détention, deux
lignes d'une même société peuvent se traiter à plusieurs dizaines de pour cent d'écart pendant des années, et le
spread est une position, et non un arbitrage.

\section{Comparer les enveloppes}

La \Cref{fig:m1:delta-one-instruments:stack} applique la
\cref{met:m1:delta-one-instruments:five} aux cinq investisseurs du chapitre pour
une année, avec les conditions illustratives du tutoriel~: un marché avec une taxe
d'achat de 50 points de base, un rendement des dividendes de 2~\% dont un
détenteur étranger perd 15~\%, un spread de \omterm{def:m1:financing:pb}{prime broker} de 50 points de base, un spread sur contrats à terme de 30, un
swap à 40 avec 95~\% du dividende, un fonds facturant 7, et un CFD de détail à
250.

\begin{omfigure}
\begin{tikzpicture}
\begin{axis}[width=11.5cm,height=5.6cm,ybar stacked,bar width=13pt,
  ylabel={coût sur un an (pb)},
  symbolic x coords={shares,ETF,future,swap,CFD},xtick={shares,ETF,future,swap,CFD},xticklabels={{actions},{ETF},{future},{swap},{CFD}},
  tick label style={font=\small},label style={font=\small},xticklabel style={font=\small\strut},
  ymin=0,ymax=330,ymajorgrids,grid style={gray!25},enlarge x limits=0.12,
  legend columns=5,legend style={font=\footnotesize,at={(0.5,-0.18)},anchor=north,draw=none,fill=none,/tikz/every even column/.append style={column sep=6pt}},
  table/col sep=comma]
\addplot[fill=omExo!50,draw=omExo] table[x=wrapper,y=trading]{figdata/markets-1/17-delta-one-instruments/stack_leveraged.csv};
\addplot[fill=omThm!70,draw=omThm] table[x=wrapper,y=tax]{figdata/markets-1/17-delta-one-instruments/stack_leveraged.csv};
\addplot[fill=black!60,draw=black] table[x=wrapper,y=running]{figdata/markets-1/17-delta-one-instruments/stack_leveraged.csv};
\addplot[fill=omDef!60,draw=omDef] table[x=wrapper,y=funding]{figdata/markets-1/17-delta-one-instruments/stack_leveraged.csv};
\addplot[fill=omDef!15,draw=omDef] table[x=wrapper,y=dividends]{figdata/markets-1/17-delta-one-instruments/stack_leveraged.csv};
\legend{négociation,taxe,frais,financement,dividendes}
\end{axis}
\end{tikzpicture}

\omcaption{Un an de la même exposition pour un investisseur étranger à effet de levier~:
140, 93, 68, 56 et 296 points de base. Un investisseur disposant de liquidités ne paie aucun financement
sur les actions ni sur le fonds, et l'ordre change~: fonds 43, swap 56, future 68,
actions 90, CFD 296. Données~: comparateur du tutoriel, avec des
termes illustratifs.}\label{fig:m1:delta-one-instruments:stack}
\end{omfigure}

\begin{omfigure}
\begin{tikzpicture}
\begin{axis}[width=11cm,height=5.2cm,
  xlabel={durée de détention (années)},ylabel={coût annuel (pb)},
  tick label style={font=\small},label style={font=\small},
  xmin=0,xmax=5,ymin=0,ymax=200,grid=major,grid style={gray!25},
  legend columns=4,legend style={font=\footnotesize,at={(0.5,-0.3)},anchor=north,draw=none,fill=none,/tikz/every even column/.append style={column sep=8pt}},
  table/col sep=comma]
\addplot[omThm,very thick] table[x=years,y=shares]{figdata/markets-1/17-delta-one-instruments/horizon.csv};
\addplot[black,thick,dashdotted] table[x=years,y=ETF]{figdata/markets-1/17-delta-one-instruments/horizon.csv};
\addplot[omDef,very thick,dashed] table[x=years,y=future]{figdata/markets-1/17-delta-one-instruments/horizon.csv};
\addplot[omExo,very thick,dotted] table[x=years,y=swap]{figdata/markets-1/17-delta-one-instruments/horizon.csv};
\legend{actions,ETF,future,swap}
\end{axis}
\end{tikzpicture}

\omcaption{Coût annualisé en fonction de la durée de détention pour un investisseur disposant de liquidités.
Les coûts initiaux (la taxe à l'achat sur les actions) sont amortis~; les coûts courants
(le financement intégré aux contrats à terme et aux swaps) ne le sont pas. Les actions dépassent le swap au-delà de
2.7 ans et le contrat à terme au-delà de 1.6~; le fonds, exonéré ici de la taxe, est
le moins cher au-delà de sept semaines. Données~: le comparateur du
tutoriel.}\label{fig:m1:delta-one-instruments:horizon}
\end{omfigure}

La comparaison est mécanique dès que les cinq nombres sont connus~; le travail
consiste à les connaître. Trois sont cotés (négociation, frais, \omterm{def:m1:delta-one-instruments:swap}{spread de financement})~; le
coût fiscal dépend du domicile et du statut de l'investisseur~; et la part des
dividendes $\eta$ d'un swap est l'estimation privée du dealer de ce que les actions
rapportent entre ses mains. Un \omterm{def:m1:delta-one-instruments:deltaone}{desk delta-one} gagne sa vie grâce aux écarts de
quelques points de base sur ces nombres d'un client à l'autre, et un investisseur
quantitatif perd un montant comparable à force de ne jamais poser la question.

\section{Tutoriel~: cinq véhicules, une même exposition}\label{tut:m1:delta-one-instruments:five}

\begin{tutorial}
\textbf{Objectif.} Calculer le coût total de cinq véhicules pour deux investisseurs
et un horizon quelconque, ainsi que la bande de conversion d'un \omterm{def:m1:delta-one-instruments:dr}{certificat de dépôt}. \textbf{État final~:}
les trois figures de données de ce chapitre.
\begin{tutsteps}
\item \textbf{Termes.} Un enregistrement par véhicule, cinq types de coûts.
\omcode{code/markets-1/17-delta-one-instruments/python/wrappers.py}{19}{25}{Termes illustratifs~: entrée, sortie, taxe à l'achat, frais courants, spread de financement, coût de roll, nombre de rolls par an, fuite de dividendes.}\label{lst:m1:delta-one-instruments:terms}
\item \textbf{La décomposition.} Les véhicules synthétiques financent la totalité du notionnel
dans leur prix, quel que soit leur détenteur.
\omcode{code/markets-1/17-delta-one-instruments/python/wrappers.py}{28}{41}{Les cinq coûts sur une période de détention.}\label{lst:m1:delta-one-instruments:breakdown}
\item \textbf{La bande de conversion du certificat.}
\omcode{code/markets-1/17-delta-one-instruments/python/wrappers.py}{72}{78}{Primes au-delà desquelles l'émission ou l'annulation de certificats devient rentable.}\label{lst:m1:delta-one-instruments:band}
\end{tutsteps}
\textbf{Pour aller plus loin.} Ajouter le côté vendeur~: un future vendu
\emph{rapporte} le financement élevé, un fonds vendu rapporte les frais et paie les
\omterm{def:m1:financing:seclending}{frais d'emprunt}. Puis rendre saisonnier le spread des contrats à terme, plus élevé lors du
roll de fin d'année, et voir pour quels horizons le future cesse d'être compétitif.
\end{tutorial}

\section{Construction~: le comparateur de véhicules}\label{bld:m1:delta-one-instruments:comparator}

\begin{build}
\bfield{Rôle} Avant que la société miniature ne prenne une exposition sur un indice ou
sur une seule action, longue ou courte, elle demande quel véhicule est le moins cher pour
son horizon et sa situation fiscale, et elle en consigne la raison.
\bfield{Interface} \texttt{WrapperTerms(\ldots)} avec un indicateur \texttt{synthetic}
et un indicateur \texttt{shortable}~; \texttt{Investor(financed,
dividend\_leak\_override)}~; \texttt{cost(wrapper, investor, years,
div\_yield\_bp, side, borrow\_fee\_bp)} renvoyant un total et ses composantes~;
\texttt{rank(\ldots)} renvoyant les cotations triées par total.
\bfield{Règles} Position longue~: taxe à l'achat, financement sur la fraction financée
(tout, pour les véhicules synthétiques), fuite de dividendes. Position courte~: pas de taxe à l'achat, le
spread du côté vendeur, plus les \omterm{def:m1:financing:seclending}{frais d'emprunt} vus dans la
\cref{bld:m1:stock-loan-and-short-selling:borrow}, les frais des fonds changent de signe, et
les véhicules qui ne peuvent pas être vendus à découvert sont exclus, avec une erreur si l'un est
demandé explicitement. Les ex æquo sont départagés par le nom.
\bfield{Tests d'acceptation} \texttt{code/firm/wrappers/tests/}~: les deux classements du
chapitre~; un détenteur exonéré d'impôt sur cinq ans~; le côté vendeur.
\bfield{Pour aller plus loin} Remplacer chaque estimation ponctuelle par un intervalle et signaler les
véhicules les moins chers quelque part dans la boîte~: en général deux, ce qui constitue la
réponse honnête.
\end{build}

\begin{omsources}
\item US Securities and Exchange Commission, \emph{Investor Bulletin:
American Depositary Receipts}, août 2012.
\item HM Revenue \& Customs, \emph{Tax when you buy shares}, gov.uk; Stamp
Taxes on Shares Manual, STSM042050 (intermediary relief).
\item Internal Revenue Service, Notice 2024-44 (section 871(m), mise en œuvre progressive).
\item A.~Madhavan, U.~Marchioni, W.~Li et D.~Y.~Du, «~Equity ETFs versus
index futures~: a comparison for fully funded investors~», \emph{Journal of
Index Investing}, 2014.
\end{omsources}

\section{Exercices}

\begin{exercise}[$\star$]\label{exo:m1:delta-one-instruments:1}
Un fonds reçoit le total return d'un indice sur 50~millions de dollars et verse
l'indice de référence, 4.2~\%, plus 40 points de base, actual/360. Sur une période de 91 jours,
l'indice rapporte 6~\% dividendes inclus. Donner les deux paiements et le montant net.
\end{exercise}

\begin{exercise}[$\star$]\label{exo:m1:delta-one-instruments:2}
Reprendre l'\cref{ex:m1:delta-one-instruments:implied} et donner le prix du
contrat à terme auquel son \omterm{def:m1:delta-one-instruments:swap}{spread de financement} serait nul, et le prix auquel
ce spread atteindrait 50 points de base.
\end{exercise}

\begin{exercise}[$\star$]\label{exo:m1:delta-one-instruments:3}
Un ADR représente 2 actions ordinaires. L'action ordinaire vaut 42.10~\euro, l'euro
vaut 1.17~\$, et l'ADR se négocie à 98.90~\$. Donner le prix de parité et la prime en
points de base.
\end{exercise}

\begin{exercise}[$\star\star$]\label{exo:m1:delta-one-instruments:4}
Avec les termes du chapitre, comparer les actions et le CFD pour un investisseur
disposant de liquidités. En dessous de quelle durée de détention le CFD est-il moins cher~? (La perte liée au dividende
est la même pour les deux.)
\end{exercise}

\begin{exercise}[$\star\star$]\label{exo:m1:delta-one-instruments:5}
Une action américaine rapporte 1.5~\%. Donner le coût annuel de la \omterm{def:m1:delta-one-instruments:wht}{retenue à la source} pour
un détenteur étranger au taux légal de 30~\% et au taux conventionnel de 15~\%. Un dealer propose un
swap qui verse 100~\% du dividende. En utilisant
l'\cref{dat:m1:delta-one-instruments:tax}, expliquer pourquoi l'offre devrait rendre le
détenteur méfiant.
\end{exercise}

\begin{exercise}[$\star\star$]\label{exo:m1:delta-one-instruments:6}
Pour l'investisseur disposant de liquidités, vérifier les deux points de croisement indiqués dans
la \cref{fig:m1:delta-one-instruments:horizon} (les actions contre le swap, le fonds
contre le contrat à terme), à partir des termes du tutoriel.
\end{exercise}

\begin{exercise}[$\star\star\star$]\label{exo:m1:delta-one-instruments:7}
\emph{Codage.} Exécuter le comparateur pour l'investisseur à effet de levier sur trois
mois. Donner les cinq totaux et leurs valeurs annualisées. Quels supports sont à égalité,
et qu'est-ce qui les départagerait en pratique~?
\end{exercise}

\begin{exercise}[$\star\star\star$]\label{exo:m1:delta-one-instruments:8}
\emph{Trouver le défaut.} La page d'un courtier affirme~: «~CFD~: pas de commission, pas de
\omterm{def:m1:asian-and-emerging-equity-markets:stamp}{droit de timbre}. Le moyen le moins cher de détenir des actions britanniques~». Un client le conserve pendant trois ans,
avec les termes du chapitre. Calculer les deux coûts et nommer ce que la page omet.
\end{exercise}

\section{Problème~: le certificat de dépôt et l'action ordinaire}

\begin{problem}[{Problème du week-end --- une société, deux tickers, un océan}]\label{pb:m1:delta-one-instruments:1}
L'action ordinaire d'une société britannique se négocie à Londres à 500~pence. Son ADR représente 4
actions ordinaires et se négocie à New York. La livre sterling vaut 1.25~\$. Le dépositaire
facture 5~cents de dollar par ADR pour l'émission ou l'annulation. Déposer des actions
britanniques dans le programme d'ADR coûte 1.5~\% de frais, tels que prévus par
l'\cref{dat:m1:delta-one-instruments:tax}. Négocier les deux jambes et la devise coûte
6~points de base au total.

\medskip\noindent\textbf{Partie I --- Parité et bande.}
\begin{enumerate}
\item Donner le prix de parité de l'ADR.
\item Donner la borne supérieure de la bande de conversion en points de base.
\item Donner la borne inférieure.
\item L'ADR se négocie à 25.30~\$. Donner sa prime. Une conversion est-elle rentable~?
\item Expliquer en une phrase pourquoi la bande est asymétrique, et ce que cela
prédit quant au signe de la prime moyenne.
\end{enumerate}

\medskip\noindent\textbf{Partie II --- Deux conversions.}
\begin{enumerate}[resume]
\item L'ADR vaut 25.50~\$. Une société émet 100\,000 ADR. Lister les flux de trésorerie
et donner le gain en dollars et en points de base.
\item L'ADR vaut 24.90~\$. La société annule 100\,000 ADR. Mêmes questions.
\item La conversion prend deux jours. La société laisse la devise non couverte~; avec
une volatilité de la livre sterling de 8~\% par an, donner l'écart-type
de l'effet de change sur deux jours, en points de base. Conclure.
\item Laquelle des deux conversions exige un emprunt de titres si la société veut
vendre la jambe qu'elle va recevoir \emph{avant} de la recevoir~?
\end{enumerate}

\medskip\noindent\textbf{Partie III --- Après la clôture de Londres.}
La place de Londres clôture à 11 h 30, heure de New York. À 15 h, l'ADR se négocie à 25.45~\$
alors que le dernier cours de Londres donne encore une parité de 25.00~\$. Le marché américain
a progressé de 1.5~\% depuis 11 h 30, et le bêta de l'action par rapport à lui est de 1.1.
\begin{enumerate}[resume]
\item Donner une estimation actualisée de la parité.
\item Donner la prime de l'ADR par rapport à la parité figée et à cette estimation.
\item Un collègue veut «~vendre la prime de 180 points de base~». Que vend-il
réellement, et contre quoi~?
\item Comment tester, sur l'historique, si le mouvement de l'après-midi de l'ADR
prédit l'ouverture suivante à Londres~?
\item Pourquoi cette prévisibilité n'est-elle pas un gain sans risque pour un trader
sur l'enchère d'ouverture à Londres~?
\end{enumerate}

\medskip\noindent\textbf{Partie IV --- Jugement.}
\begin{enumerate}[resume]
\item Pour un investisseur américain qui souhaite détenir cette société pendant cinq ans, lister
ce que l'ADR coûte et que l'action ordinaire ne coûterait pas, et l'inverse.
\item La société transfère sa cotation principale à New York et l'ADR est
remplacé par l'action elle-même. Que devient la bande~?
\item Deux catégories d'une même société se négocient à 30~\% d'écart dans deux pays
entre lesquels la conversion est interdite. De quelle nature est l'opération «~acheter le
moins cher, vendre le plus cher~»~?
\item Qui sont les opérateurs naturels des conversions de la partie II~?
\item Énoncer le \emph{résultat nommé}~: la bande de conversion, en points de base de
la parité.
\item En une phrase~: que vend un \omterm{def:m1:delta-one-instruments:deltaone}{desk delta-one}~?
\end{enumerate}
\end{problem}

\section{Questions d'entretien}

\begin{interviewq}[$\star$ \iqroles{trader, banque, chercheur}]\label{iq:m1:delta-one-instruments:1}
Que signifie «~\omterm{def:m1:delta-one-instruments:deltaone}{delta-one}~», et que fait un \omterm{def:m1:delta-one-instruments:deltaone}{desk delta-one}~?
\end{interviewq}

\begin{interviewq}[$\star$ \iqroles{trader, banque}]\label{iq:m1:delta-one-instruments:2}
Un client peut acheter un fonds indiciel ou se mettre acheteur de contrats à terme sur indice. Comment
comparez-vous les deux pour elle~?
\end{interviewq}

\begin{interviewq}[$\star\star$ \iqroles{trader, chercheur}]\label{iq:m1:delta-one-instruments:3}
Les contrats à terme sur indice se traitent «~chers~»~: le taux de financement implicite
dépasse de 80 points de base l'indice de référence. Qui en pâtit, qui en profite, et quelle
transaction capte le spread~?
\end{interviewq}

\begin{interviewq}[$\star\star$ \iqroles{banque, trader}]\label{iq:m1:delta-one-instruments:4}
Vous êtes dealer et cotez un \omterm{def:m1:financing:trs}{total return swap} sur une action européenne à un
\omterm{def:m1:what-a-trading-firm-does:hedge-fund}{hedge fund}. Qu'est-ce qui entre dans votre \omterm{def:m1:delta-one-instruments:swap}{spread de financement} et dans votre pourcentage de dividende~?
\end{interviewq}

\begin{interviewq}[$\star\star$ \iqroles{trader, chercheur}]\label{iq:m1:delta-one-instruments:5}
Un ADR se traite à 1~\% au-dessus de sa parité. S'agit-il d'un arbitrage~? Que vérifiez-vous~?
\end{interviewq}

\begin{interviewq}[$\star\star\star$ \iqroles{chercheur, trader}]\label{iq:m1:delta-one-instruments:6}
Votre backtest long-short sur actions suppose que vous détenez des actions. La société
la mettrait en œuvre sur swap. Qu'est-ce qui change dans le backtest~?
\end{interviewq}
